\section{Conditional composition into a material}\label{sec:conditional}
We isolate both assumptions needed to turn a referee cell into a
translation-invariant material. Neither follows merely from small
positive cell energy or from a general circuit-to-Hamiltonian citation.

\begin{definition}[Local selection hypothesis]\label{def:selection}
For every classical configuration, intact cells support disjoint
$K_h$ and all remaining quantum terms are nonnegative. The area
decomposition into wall segments has indicators $v_s$ for divider
corners and $i_s\leq v_s$ for intact-cell corners, at pitch $p_s$.
Clean matching runs have constant pitch and deficit greater than $-1$,
with zero deficit for a closed run. Open and bad runs can be charged
at most six times to each wall violation. Missing intact-cell corners
can be charged at most six times to a wall violation and once to a
tableau or flag violation. The rules specify these assignments locally.
\end{definition}

\begin{lemma}[Selection from the local hypothesis]\label{lem:positive-selection}
Assume Definition~\ref{def:selection} and
$0\leq e(h)\leq\min(1/64,ch^{-3})$. Perfect stripes are selected
by Lemma~\ref{lem:shifted}; at fixed small off-tie $g$ every
nonminimizing block has eventual separation at least $7/16$.
\end{lemma}
\begin{proof}
Write $e_s=e(p_s-1)$, $e_{\max}=1/64$ and
$f^*_e=\min_pF_g(p)<0$. The exact area bookkeeping and intact-cell
lower bound give energy excess at least
\begin{align*}
 &\sum_s p_s(F_g(p_s)-f^*_e)
 +\sum_s(r+e_s)(v_s-1/p_s)-\sum_s e_s(v_s-i_s)\\
 &\hspace{20mm}+V_w+V_t+\operatorname{Area}_{W\text{-free}}(-f^*_e).
\end{align*}
The run hypothesis bounds the second term below by
$-6(r+e_{\max})V_w$. The missing-corner hypothesis bounds the third
by $-6e_{\max}V_w-e_{\max}V_t$. Thus the defect coefficients are
at least $1-6r-12e_{\max}=7/16$ and
$1-e_{\max}=63/64$. Clean wrong stripes have positive density
separation; at sufficiently large commensurate factorial volume their
energy excess exceeds $7/16$, as does the vacuum excess.
\end{proof}
This proves the implication from the local hypothesis. The hypothesis
itself is an additional combinatorial rule-level requirement for
positive-energy cells.

\begin{definition}[Referee layout hypothesis]\label{def:layout}
A fixed finite alphabet and rational algebraic onsite/bond rule set
must implement, inside each selected $h\times h$ cell, the following.
The depth is $N=2\lfloor(c/2)\log_2h\rfloor$, the strategies are those
of Proposition~\ref{prop:circuits}, and the padded clocks have
$h^{3/2}\leq T\leq h^C$. There is an exact invariant Cartesian legal
propagation space with product rewinding and either the corner
initialization in \eqref{eq:cell} or the commuting row penalties used
in \eqref{eq:rowconstant}. The loss and question predicates have
fixed-size interface support; initialization, transport and folds have
specified nearest-neighbour rules. A local isometric support description
gives $\log_2D_A=O(N^2+\log T)$ for low-energy cell vectors and
their reference states. Unused sites are frozen. The generated
classical gate labels fit the cell's time and space budget. All
nonlegal sectors have the claimed quantitative separation and the
factorized projection/product extension required for any asserted
physical-state extraction. The effective extraction constant $C'_N$
is polynomial in $N$ and the reference distance is
$\operatorname{poly}(N)\sqrt{M\lambda}$.
\end{definition}

For an assertion about Gibbs states one must additionally have a genuine
bipartite tensor-product state space and a proved bounded cut norm.
An abstract globally conditioned legal subspace alone does not supply
these thermal hypotheses.

\begin{theorem}[Conditional halting-game cell coefficient]\label{thm:conditional}
Assume Definitions~\ref{def:selection} and \ref{def:layout}, with
$c=3/4$, the admitted dimension-matched halting-game strategies, and
the specified competing/frozen margins. In every selected cell, every
pure ground vector satisfies
\begin{equation}\label{eq:conditional}
 S_A=4\alpha_KN+o(N)
    =3\alpha_K\log_2(1/g)+o(\log(1/g)).
\end{equation}
For any marginal $\rho_c$ supported on that cell ground space,
\begin{equation}\label{eq:marginal}
 \ED(\rho_c)=4\alpha_KN+o(N),\qquad
 S((\rho_c)_A)=4\alpha_KN+o(N).
\end{equation}
The coefficient $3\alpha_K=3(3+\theta_K)/4$ is Turing-equivalent to
$K$. This is a conditional cell theorem; no unconditional
translation-invariant implementation is asserted.
\end{theorem}
\begin{proof}
The initialized test strategy has target loss budget
$\eps_N=C_1 16^{-N}$, so $Me(h)\leq\eps_N$ and
Lemma~\ref{lem:padding} gives $e(h)\leq h^{-3}$.
Freeze the finitely many cells below the fixed threshold needed for
$e(h)\leq1/64$. Lemma~\ref{lem:positive-selection} and the assumed
margins then give selected-block transfer.
For every pure cell ground vector, extraction yields
\[
 S_A\geq d_{G_K}^{\to}(C'_N\eps_N)=4\alpha_KN+o(N),
\]
because $\log_2C'_N=O(\log N)$ and
$\log_2(1/\eps_N)=4N+O(1)$. The actual optimal loss is only upper
bounded by this target budget and is not identified with it.
The normalized reference distance is
$\operatorname{poly}(N)4^{-N}$ and the assumed support logarithm
is polynomial. Proposition~\ref{prop:factorization}, or its row
version, gives
\[
 S_A\leq\log_2(c_0d_N)+\log_2|X|+o(1)=4\alpha_KN+o(N).
\]
This proves the pure squeeze.

For the mixed marginal, the loss operator inequality gives the same
lower usable price. Uniformity of the projection estimate over the
ground space bounds its weight outside the reference subspace by
$\operatorname{poly}(N)16^{-N}$. Project and normalize the density
matrix. The gentle projection bound gives trace distance
$\operatorname{poly}(N)4^{-N}$; its reference Alice support has
dimension at most $|X|c_0d_N$. Entropy continuity with the same actual
support bound gives
\[
 4\alpha_KN+o(N)\leq\ED(\rho_c)\leq S((\rho_c)_A)
                         \leq4\alpha_KN+o(N).
\]
The middle inequality is the standard upper bound of distillation by
local entropy. The last entropy need not itself be within-cell
entanglement when $\rho_c$ is mixed. This argument does not imply
entropy additivity for arbitrary unions of cells.

Lemma~\ref{lem:shifted} gives
$\log_2h=\log_2(1/g)+O(1)$; the even depth rounding changes $N$
by at most two, so \eqref{eq:conditional} follows. An oracle for
$K$ computes $\theta_K$ by finite sums and the geometric tail.
Conversely subtract known contributions below $e$ and put
$w_e=4^{-(e+1)}$. If $e\notin K$ the residual is at most $w_e/3$;
if $e\in K$ it is at least $w_e$. Accuracy $w_e/6$ separates the
two cases, inductively computing $K$. Rational affine rescaling
preserves this Turing degree.
\end{proof}

The resource inequalities motivate the hypothesis but do not prove its
rules. For $N=(3/4)\log_2h+O(1)$ and small fixed $\delta<1/3$,
Proposition~\ref{prop:circuits} gives generation time $o(h)$ after
absorbing its polynomial factor. Its $O(N^2)$ work qubits and padded
length $h^{3/2}$ use $h^{3/2}\operatorname{poly}(\log h)=o(h^2)$
time-space area. Exact legal propagation after folding, local endpoint
access, distributed initialization, and label packing still require
an actual rule table. The complete layout is a plausible open
construction problem, stated once here. We do not extend this
conditional result to every game with a limiting price coefficient.
